% Created 2026-03-20 sex 11:59
% Intended LaTeX compiler: lualatex
\documentclass[12pt]{article}
\usepackage[a4paper,left=3cm,top=3cm,right=2cm,bottom=2cm]{geometry}
\setcounter{secnumdepth}{4}
\setcounter{tocdepth}{2}
\usepackage{graphicx}
\usepackage{longtable}
\usepackage{wrapfig}
\usepackage{rotating}
\usepackage[normalem]{ulem}
\usepackage{amsmath}
\usepackage{amssymb}
\usepackage{capt-of}
\usepackage{hyperref}
\usepackage{fgv-header}
\usepackage{float}
\usepackage{svg}
\author{Gustavo M. Mendes de Tarso}
\date{\today}
\title{}
\hypersetup{
 pdfauthor={Gustavo M. Mendes de Tarso},
 pdftitle={},
 pdfkeywords={},
 pdfsubject={},
 pdfcreator={Emacs 28.2 (Org mode 9.5.5)}, 
 pdflang={English}}
\begin{document}

\begingroup
\linespread{1}\selectfont
\begin{tcolorbox}[title=Ficha Técnica,
  colback=gray!5,colframe=gray!40,boxrule=0.4pt,sharp corners]
\begin{tblr}{rowsep=1pt,stretch=1, rows={t}, colspec={Q[l,2.8cm] X[l]}}
\textbf{Curso}                   & Mestrado de Políticas Públicas e Governo \\
\textbf{Turma}                   & T-01 \\
\textbf{Pólo}                    & Brasília \\
\textbf{Disciplina}              & Teorias de Administração Pública \\
\textbf{Professor}               & Bernardo Buta \\
\textbf{Trabalho}                & Busca de literatura em dupla \\
\textbf{Aluno(s)}                & Gustavo M. Mendes de Tarso \\
\textbf{Título do trabalho}      & Busca, seleção e síntese de artigo de revisão sobre stakeholders a partir do Semantic Scholar com aplicação do PRISMA 2020 \\
\end{tblr}
\end{tcolorbox}
\endgroup

\FGVNewPage

\section{Introdução}
\label{sec:org304eb1a}

Este trabalho apresenta, em formato compatível com o \textbf{PRISMA 2020}, o processo de \textbf{identificação, triagem, elegibilidade e inclusão} de um \textbf{artigo de revisão} sobre \textbf{stakeholders}, tomando como base uma \textbf{busca focalizada no Semantic Scholar}. O objetivo não foi produzir uma revisão sistemática exaustiva de toda a literatura internacional, mas demonstrar de modo transparente um percurso de busca acadêmica, seleção e justificativa de escolha, tal como solicitado na atividade.

O tema foi definido a partir do termo \textbf{stakeholders}, por ser um conceito central em debates contemporâneos sobre governança, gestão, engajamento, participação e formulação estratégica. Como o enunciado da aula enfatiza a \textbf{literatura recente} e a necessidade de selecionar \textbf{artigos de revisão}, a busca privilegiou estudos publicados nos últimos anos e com expressões como \textbf{systematic review}, \textbf{literature review} ou equivalentes no título e/ou nos metadados recuperados pelo Semantic Scholar.

O artigo selecionado ao final do processo foi \textbf{“Exploring the antecedents and consequences of firm-stakeholder engagement process: A systematic review of literature”}, de \textbf{Avinash Pratap Singh} e \textbf{Zillur Rahman} (2022). A escolha se justifica porque, entre os resultados recentes encontrados, este foi o texto com \textbf{maior aderência geral ao problema dos stakeholders}, sem se restringir a um único setor específico, como saúde, proteína alternativa, ensino superior ou logística reversa. Além disso, o estudo oferece um desenho de revisão mais robusto, articulando \textbf{antecedentes, moderadores e consequências} do processo de \textbf{stakeholder engagement}.

Também é importante registrar uma distinção metodológica. Neste trabalho, o \textbf{PRISMA 2020} foi utilizado para organizar o \textbf{fluxo da minha busca e seleção no Semantic Scholar}. Já o artigo selecionado adota, internamente, um desenho próprio de revisão sistemática, baseado em \textbf{Tranfield, Denyer e Smart (2003)}.

\section{Estratégia de busca}
\label{sec:org466a398}

\subsection{Pergunta orientadora}
\label{sec:orga9df4ea}

\begin{quote}
Qual artigo de revisão recente, localizado via Semantic Scholar, melhor sintetiza a literatura sobre stakeholders de forma ampla, sem restringir o conceito a um único domínio empírico altamente setorial?
\end{quote}

\subsection{Base e lógica de busca}
\label{sec:orgc330d1d}

A busca foi realizada com foco no \textbf{Semantic Scholar}, usando termos que combinassem o núcleo conceitual do tema com marcadores de revisão de literatura. Como o mecanismo da plataforma opera por relevância semântica, a estratégia foi organizada como uma \textbf{busca focalizada e didática}, seguida de triagem manual por título, resumo disponível, metadados e, quando possível, texto completo.

\subsection{String de trabalho}
\label{sec:org3ef1489}

\begin{quote}
("stakeholder" OR "stakeholders" OR "stakeholder engagement") AND ("systematic review" OR "literature review" OR review)
\end{quote}

\subsection{Critérios de inclusão}
\label{sec:orgf274cdc}

\begin{itemize}
\item artigo de revisão ou revisão sistemática;
\item aderência explícita ao tema \textbf{stakeholders};
\item publicação recente, privilegiando o recorte \textbf{2022-2025};
\item disponibilidade de metadados suficientes para avaliação;
\item preferência por textos com escopo \textbf{geral} ou \textbf{transversal}, e não estritamente setorial.
\end{itemize}

\subsection{Critérios de exclusão}
\label{sec:org084c53e}

\begin{itemize}
\item estudos excessivamente setoriais, quando o conceito de stakeholders aparecia restrito a um único domínio empírico;
\item textos cujo foco principal não era propriamente a literatura sobre stakeholders, mas um campo aplicado muito delimitado;
\item resultados menos aderentes ao objetivo da atividade, mesmo quando formalmente classificados como revisão.
\end{itemize}

\section{PRISMA 2020}
\label{sec:org21d85f4}

\subsection{Observação metodológica}
\label{sec:orgfac8bfc}

O fluxo abaixo corresponde à \textbf{minha busca focalizada no Semantic Scholar}, e não ao fluxo interno de seleção do artigo escolhido. Em outras palavras, a matriz PRISMA 2020 aqui apresentada serve para mostrar \textbf{quantos registros foram encontrados, quantos foram excluídos e por quê}, até a definição do artigo final da atividade.

\subsection{Fluxo textual dos estudos}
\label{sec:org33553ec}

\subsubsection{Identificação}
\label{sec:orgdb57fb8}
\begin{itemize}
\item Registros identificados por busca focalizada no \textbf{Semantic Scholar}: \textbf{n = 6}
\item Registros removidos antes da triagem: \textbf{n = 0}
\begin{itemize}
\item Duplicatas removidas: \textbf{n = 0}
\item Remoção por outros motivos: \textbf{n = 0}
\end{itemize}
\end{itemize}

\subsubsection{Triagem}
\label{sec:org0382415}
\begin{itemize}
\item Registros triados por título, resumo e metadados: \textbf{n = 6}
\item Registros excluídos na triagem: \textbf{n = 2}
\begin{itemize}
\item Revisão com foco em \textbf{saúde materno-infantil} em países de baixa e média renda: \textbf{n = 1}
\item Revisão com foco em \textbf{logística reversa no comércio eletrônico}: \textbf{n = 1}
\end{itemize}
\end{itemize}

\subsubsection{Elegibilidade}
\label{sec:org6d747b1}
\begin{itemize}
\item Relatórios buscados para recuperação em texto completo: \textbf{n = 4}
\item Relatórios não recuperados: \textbf{n = 0}
\item Relatórios avaliados para elegibilidade: \textbf{n = 4}
\item Relatórios excluídos após leitura do texto completo: \textbf{n = 3}
\begin{itemize}
\item Revisão centrada em \textbf{proteínas alternativas}: \textbf{n = 1}
\item Revisão centrada em \textbf{governança sustentável de cadeias de suprimento}: \textbf{n = 1}
\item Revisão centrada em \textbf{ensino superior e IA generativa}: \textbf{n = 1}
\end{itemize}
\end{itemize}

\subsubsection{Inclusão}
\label{sec:org5bf9209}
\begin{itemize}
\item Estudos incluídos na síntese qualitativa final: \textbf{n = 1}
\end{itemize}

\subsection{Tabela-resumo do fluxo PRISMA 2020}
\label{sec:org5809864}

\begin{table}[htbp]
\centering
\scriptsize
\setlength{\tabcolsep}{4pt}
\renewcommand{\arraystretch}{1.15}
\begin{tabularx}{\textwidth}{>{\raggedright\arraybackslash}X >{\raggedleft\arraybackslash}m{0.18\textwidth}}
\toprule
Etapa PRISMA 2020 & n \\
\midrule
Registros identificados & 6 \\
Registros removidos antes da triagem & 0 \\
Registros triados & 6 \\
Registros excluídos na triagem & 2 \\
Relatórios buscados para recuperação & 4 \\
Relatórios não recuperados & 0 \\
Relatórios avaliados para elegibilidade & 4 \\
Relatórios excluídos após leitura completa & 3 \\
Estudos incluídos na síntese qualitativa & 1 \\
\bottomrule
\end{tabularx}
\normalsize
\end{table}

\subsection{Estudos classificados no fluxo}
\label{sec:orgbd1f9b0}

\begin{table}[htbp]
\centering
\scriptsize
\setlength{\tabcolsep}{3pt}
\renewcommand{\arraystretch}{1.15}
\begin{tabularx}{\textwidth}{>{\raggedright\arraybackslash}X >{\raggedright\arraybackslash}m{0.18\textwidth} >{\raggedright\arraybackslash}m{0.14\textwidth} >{\raggedright\arraybackslash}m{0.22\textwidth}}
\toprule
Estudo & Situação no fluxo & Tipo & Motivo \\
\midrule
Singh, A. P.; Rahman, Z. (2022), \textit{Exploring the antecedents and consequences of firm-stakeholder engagement process: A systematic review of literature} & Incluído & Revisão sistemática & Texto mais amplo e mais aderente ao tema geral dos stakeholders \\
Amato, M. et al. (2023), \textit{Stakeholder Beliefs about Alternative Proteins: A Systematic Review} & Excluído na elegibilidade & Revisão sistemática & Recorte substantivo restrito ao debate sobre proteínas alternativas \\
Wadhwa, J. et al. (2024), \textit{Comprehensive literature review on stakeholder prospects, responsibilities, and quality performance measurement in sustainable supply chain governance} & Excluído na elegibilidade & Revisão de literatura & Foco concentrado em governança sustentável da cadeia de suprimentos \\
Yang, J.; Qiu, H.; Yu, W. (2025), \textit{Surging currents: a systematic review of the literature on dynamic stakeholder engagements in higher education in the generative artificial intelligence era} & Excluído na elegibilidade & Revisão sistemática & Enfoque muito específico em ensino superior e IA generativa \\
Singh, D. R. et al. (2024), \textit{Understanding service users and other stakeholders' engagement in maternal and newborn health services research: A systematic review of evidence from low- and middle-income countries} & Excluído na triagem & Revisão sistemática & Escopo temático muito específico em saúde materno-infantil \\
González Romero, I. et al. (2025), \textit{The perspective of key stakeholders on the sustainable design of reverse logistics in e-commerce: an interpretive sensemaking review} & Excluído na triagem & Revisão interpretativa & Recorte empírico delimitado à logística reversa no comércio eletrônico \\
\bottomrule
\end{tabularx}
\normalsize
\end{table}

\clearpage

\subsection{Representação esquemática do fluxo PRISMA 2020}
\label{sec:orgdd2a3cb}

Nesta versão, o fluxograma foi mantido como \textbf{imagem vetorial externa em SVG}, referenciada no próprio arquivo \texttt{.org}. Esse pipeline foi adotado para preservar maior fidelidade visual ao padrão oficial do \textbf{PRISMA 2020}, evitando as limitações de renderização típicas de geradores automáticos de diagrama.

\begin{figure}[H]
\centering
\includesvg[width=0.95\textwidth]{prisma_stakeholders_semantic_scholar}
\caption{Fluxograma PRISMA 2020 da seleção do artigo de revisão sobre stakeholders.}
\label{fig:prisma2020_stakeholders_fluxo}
\end{figure}

\section{Resumo do artigo selecionado}
\label{sec:orgc9a0a3d}

\subsection{Referência do estudo}
\label{sec:org0e0df43}

Singh, Avinash Pratap; Rahman, Zillur. \textbf{Exploring the antecedents and consequences of firm-stakeholder engagement process: A systematic review of literature}. \emph{Corporate Governance and Sustainability Review}, v. 6, n. 3, p. 28-39, 2022.

\subsection{Problema e objetivo de pesquisa}
\label{sec:org2e28fe5}

O artigo parte do diagnóstico de que a literatura sobre \textbf{firm-stakeholder engagement} cresceu consideravelmente, mas permanece \textbf{fragmentada}. Em muitos casos, os estudos tratam o engajamento apenas como uma atividade empresarial ou enfatizam apenas um lado da relação, seja o contexto da firma, seja o contexto dos stakeholders. Com isso, a compreensão do processo de engajamento tende a ficar dispersa, dificultando uma visão integrada de seus antecedentes e resultados.

Diante desse problema, o objetivo do artigo é \textbf{consolidar a literatura} sobre o processo de engajamento firma-stakeholder, identificando seus \textbf{antecedentes, moderadores e consequências}. Além disso, o estudo busca examinar se a literatura sustenta a existência de um \textbf{continuum} entre dimensões \textbf{instrumentais}, \textbf{integradas} e \textbf{normativas} do relacionamento com stakeholders.

\subsection{Argumento central}
\label{sec:org90ce413}

O argumento central do artigo é que o \textbf{stakeholder engagement} deve ser compreendido como um \textbf{processo relacional e contínuo}, e não como uma ação pontual de comunicação ou mera gestão tática de interesses. Nessa perspectiva, o engajamento passa a ser entendido como um mecanismo que conecta estratégia, governança, inovação, responsabilidade social corporativa e criação de valor compartilhado.

Os autores sustentam que os resultados mais relevantes do engajamento surgem quando a organização não trata stakeholders apenas como pressões externas a serem acomodadas, mas como participantes efetivos de processos de aprendizagem, cooperação, legitimidade e geração de benefícios mútuos. Assim, a contribuição do artigo é oferecer uma visão integrada do engajamento firma-stakeholder como processo cíclico e multidimensional.

\subsection{Desenho de pesquisa}
\label{sec:org84c156c}

O estudo selecionado é uma \textbf{revisão sistemática da literatura}. Os autores afirmam que adotaram como base o framework metodológico de \textbf{Tranfield, Denyer e Smart (2003)}, estruturado em três etapas: \textbf{planejamento, condução e análise}. Na etapa de condução, a busca inicial em três bases eletrônicas resultou em \textbf{1.280 artigos}. Em seguida, \textbf{687 duplicatas} foram removidas com apoio do \textbf{Zotero}, restando \textbf{593 estudos} para triagem por título e resumo. Após essa triagem, \textbf{418 artigos} foram excluídos, e \textbf{175 trabalhos} permaneceram para análise descritiva e temática.

Com esse corpus final, os autores codificaram os estudos por ano, país e método de pesquisa, além de procederem a uma análise temática dos antecedentes, moderadores e consequências do engajamento. Portanto, trata-se de uma revisão sistemática estruturada, com base quantitativa de seleção e posterior síntese qualitativa.

\subsection{Principais achados}
\label{sec:orgd1460cd}

Os resultados mostram que o engajamento com stakeholders aparece na literatura como elemento relevante em áreas como \textbf{estratégia corporativa, governança, inovação, operações e responsabilidade social corporativa}. A revisão indica que, em muitos estudos, o foco predominante ainda é \textbf{instrumental}, isto é, orientado por objetivos estratégicos da firma. Contudo, os autores também identificam espaço importante para dimensões \textbf{normativas} e \textbf{integradas}, sobretudo quando a literatura discute governança corporativa, criação de valor compartilhado e legitimidade organizacional.

Entre os antecedentes mais recorrentes, destacam-se fatores ligados ao \textbf{contexto da firma}, ao \textbf{contexto dos stakeholders} e ao \textbf{contexto institucional}. Entre as consequências, aparecem resultados para a própria organização, para os stakeholders e para uma esfera compartilhada, como \textbf{aprendizagem coletiva, reputação, legitimidade, cooperação, inovação e criação de valor conjunto}. O texto conclui que o engajamento tende a ser mais robusto quando incorporado ao núcleo das atividades organizacionais, e não tratado como mecanismo periférico.

\subsection{Justificativa da seleção final}
\label{sec:org950c39b}

O artigo de Singh e Rahman foi escolhido porque foi o estudo que melhor atendeu, simultaneamente, a quatro requisitos da atividade: \textbf{recência}, \textbf{caráter de revisão}, \textbf{centralidade do tema stakeholders} e \textbf{escopo amplo}. Os demais textos recuperados no fluxo PRISMA eram úteis, mas apareciam fortemente condicionados por um domínio empírico específico, como saúde, cadeia de suprimentos, proteínas alternativas, ensino superior ou logística reversa.

Já o artigo selecionado oferece uma síntese mais geral do debate sobre stakeholders, especialmente no eixo do \textbf{engajamento firma-stakeholder}, o que o torna mais apropriado para uma atividade de sala voltada à discussão conceitual e metodológica de um tema transversal.

\section{Texto corrido para entrega}
\label{sec:org98b1cc2}

A busca de literatura foi realizada no \textbf{Semantic Scholar}, com foco em \textbf{artigos de revisão recentes} sobre \textbf{stakeholders}. Para isso, foi utilizada uma estratégia de busca focalizada com os termos \textbf{stakeholder}, \textbf{stakeholders}, \textbf{stakeholder engagement}, \textbf{systematic review} e \textbf{literature review}. Em seguida, os resultados foram organizados segundo a lógica do \textbf{PRISMA 2020}, distinguindo as etapas de identificação, triagem, elegibilidade e inclusão. A busca inicial levou à seleção de \textbf{seis registros} diretamente aderentes ao tema. Não houve duplicatas. Após a leitura de títulos, resumos e metadados, \textbf{dois estudos} foram excluídos ainda na triagem por apresentarem escopo excessivamente setorial, um em saúde materno-infantil e outro em logística reversa no comércio eletrônico. Restaram \textbf{quatro relatórios} para recuperação e avaliação em texto completo. Desses, \textbf{três} foram excluídos na etapa de elegibilidade por manterem foco muito delimitado em proteínas alternativas, cadeia de suprimentos sustentável e ensino superior na era da IA generativa. Assim, apenas \textbf{um estudo} foi incluído na síntese qualitativa final: o artigo de \textbf{Singh e Rahman (2022)}, intitulado \textbf{Exploring the antecedents and consequences of firm-stakeholder engagement process: A systematic review of literature}. O texto foi selecionado porque apresenta o tratamento mais amplo e teoricamente consistente do tema, sem restringir stakeholders a um único campo aplicado. O artigo parte do problema de que a literatura sobre \textbf{firm-stakeholder engagement} está fragmentada e busca consolidar os \textbf{antecedentes, moderadores e consequências} desse processo. Metodologicamente, trata-se de uma revisão sistemática baseada em \textbf{Tranfield, Denyer e Smart (2003)}, com busca inicial de \textbf{1.280 artigos}, remoção de \textbf{687 duplicatas}, triagem de \textbf{593 estudos} e análise final de \textbf{175 publicações}. O argumento central é que o engajamento com stakeholders deve ser compreendido como um \textbf{processo relacional, contínuo e multidimensional}, que articula estratégia, governança, inovação e criação de valor compartilhado. Por isso, o artigo se mostrou o mais adequado para a atividade, tanto pela aderência conceitual quanto pela robustez do desenho de revisão.

\section{Referências}
\label{sec:orgeaec8a5}

\begin{itemize}
\item PAGE, Matthew J. et al. \textbf{The PRISMA 2020 statement: an updated guideline for reporting systematic reviews}. BMJ, 2021.
\item SINGH, Avinash Pratap; RAHMAN, Zillur. \textbf{Exploring the antecedents and consequences of firm-stakeholder engagement process: A systematic review of literature}. Corporate Governance and Sustainability Review, v. 6, n. 3, p. 28-39, 2022.
\end{itemize}
\end{document}